\section{Implicit constraints}\label{sec:implicit_constraints}

This appendix complements Section~\ref{sec:constraints} with inequalities that hold for the model but are not stated there explicitly.

\paragraph{Closed nodes send no flow.}
\begin{align}
    b_{st} & \le \overline{W}_s\,\alpha_s, \quad \forall s \in \mathcal{S},\ t \in \mathcal{T}, \label{eq:closed_source} \\
    c_{tz} & \le \overline{V}_t\,\beta_t, \quad \forall t \in \mathcal{T},\ z \in \mathcal{Z}. \label{eq:closed_plant}
\end{align}
By \eqref{eq:nonnegativity} and \eqref{eq:source_flow_conservation}, $b_{st} \le \sum_{t' \in \mathcal{T}} b_{st'} = a_s$, and $a_s \le \overline{W}_s\alpha_s$ by \eqref{eq:source_capacity}.
Likewise, \eqref{eq:nonnegativity} and \eqref{eq:plant_flow_conservation} give $c_{tz} \le \sum_{z' \in \mathcal{Z}} c_{tz'} = \sum_{s \in \mathcal{S}} b_{st}$, which is at most $\overline{V}_t\beta_t$ by \eqref{eq:plant_capacity}.
Hence a closed source or plant never sends water outwards.

\paragraph{Total supply covers total demand.}
Summing \eqref{eq:demand} over all zones and applying \eqref{eq:plant_flow_conservation}, \eqref{eq:source_flow_conservation} and \eqref{eq:source_capacity} gives
\begin{equation}\label{eq:total_supply}
    \sum_{s \in \mathcal{S}} \overline{W}_s\,\alpha_s
    \ \ge\ \sum_{s \in \mathcal{S}} a_s
    \ =\ \sum_{t \in \mathcal{T}} \sum_{z \in \mathcal{Z}} c_{tz}
    \ \ge\ \sum_{z \in \mathcal{Z}} D_z.
\end{equation}
Whatever the routing, the active sources must have enough combined capacity to cover total demand, so no set of sources whose maximum withdrawals sum to less than $\sum_z D_z$ can be the active set.

\paragraph{Blend bounds implied by outlet quality.}
Combining \eqref{eq:outlet_water_quality} with \eqref{eq:removal_process_initial} gives, for every $t \in \mathcal{T}$ and $p \in \mathcal{P}$:
\begin{align}
    \underline{Q}_p \sum_{s \in \mathcal{S}} b_{st}
     & \le \sum_{s \in \mathcal{S}} Q_{sp}\, b_{st}, \label{eq:implied_blend_lower}
\end{align}
This means the incapability of $t \in \mathcal T$ to correct a blend that is below a lower limit, assuming only removal process is implemented; thereby, every lower regulatory limit must be met by source selection and blending, $\forall p \in \mathcal P$.
Three cases follow.
\begin{itemize}
    \item If $\underline{Q}_p = 0$, \eqref{eq:implied_blend_lower} never binds, otherwise no water is drawn.
    \item If $\underline{Q}_p > 0$, as for alkalinity, \eqref{eq:implied_blend_lower} binds regardless of the processes installed at the plant.
    \item If no process acts on $p$ ($H_{trp} = 0$ for all $r \in \mathcal{R}_t$, or equivalently $e_{trp} = 0$, as for the transformed value \eqref{eq:ph_source_transformation}), then $\ell^{n_t}_{tp} = \ell^{0}_{tp}$ and
          \begin{equation*}
              \underline{Q}_p \sum_{s \in \mathcal{S}} b_{st}
              \le \sum_{s \in \mathcal{S}} Q_{sp}\, b_{st}
              \le \overline{Q}_p \sum_{s \in \mathcal{S}} b_{st},
          \end{equation*}
          so both limits must be met by blending alone.
\end{itemize}